\section{Coverage and Verification Guarantees}\label{sec:theory}

This section separates three claims that are often conflated: a finite
proposal cannot cover arbitrary targets; it can cover targets satisfying a
structural transfer condition; and a recommendation can be verified without
assuming the search model is correct.  Complete machine-checked statements are
mapped to the implementation in Appendix \ref{app:lean}.

\subsection{Why an Explicit Transfer Condition Is Necessary}

\begin{theorem}[Finite-atlas no-free-lunch]\label{thm:nfl}
Let $\cX$ be finite and let $\cA\subsetneq\cX$ be any proper proposal atlas
frozen before target labels.  There exists a nonempty feasible set
$\cF\subseteq\cX$ such that $\cA\cap\cF=\varnothing$.
\end{theorem}

\begin{proof}
Choose $x_0\in\cX\setminus\cA$ and set $\cF=\{x_0\}$.  Then $\cF$ is
nonempty and is missed by $\cA$.
\end{proof}

Theorem \ref{thm:nfl} is elementary but decisive.  Neither source learning nor
a nonlinear encoder can guarantee that ten target-label-free proposals hit an
arbitrary feasible set.  A positive result must identify what transfers.

\subsection{Aligned Geometric Atlas Coverage}

Let $\etaStar$ be an ideal transferable coordinate and $\etaHat$ the coordinate
learned from source information.  Let $\cS$ be the finite source-support
library and $\atlas\subseteq\cS$ the selected atlas.  Distances below are in
the common coordinate space.

\begin{assumption}[Atlas coverage]\label{ass:atlas-cover}
For every $s\in\cS$, there is $a\in\atlas$ such that
$\|\etaHat(s)-\etaHat(a)\|\le r_{\mathrm{cover}}$, with
$|\atlas|\le n_0$.
\end{assumption}

\begin{assumption}[Coordinate approximation]\label{ass:coordinate}
For each relevant policy $x$,
$\|\etaHat(x)-\etaStar(x)\|\le\epsilon_\eta$.
\end{assumption}

\begin{assumption}[Source-target support]\label{ass:support}
There exists a target-safe center $x_c$ and a source-support policy
$s_c\in\cS$ such that
$\|\etaStar(x_c)-\etaStar(s_c)\|\le\Delta_{\mathrm{task}}$.
\end{assumption}

\begin{assumption}[One-sided margin regularity and safe depth]
\label{ass:margin}
The target chance margin satisfies
\begin{equation}\label{eq:lipschitz}
 m_q(x)\le m_q(y)+L_m\|\etaHat(x)-\etaHat(y)\|
\end{equation}
for relevant $x,y$, and $m_q(x_c)+\gamma\le0$ for some $\gamma>0$.
\end{assumption}

\begin{theorem}[Aligned finite-atlas coverage]\label{thm:coverage}
Under Assumptions \ref{ass:atlas-cover}--\ref{ass:margin}, if
\begin{equation}\label{eq:coverage-condition}
 L_m\bigl(r_{\mathrm{cover}}+\Delta_{\mathrm{task}}
       +2\epsilon_\eta\bigr)\le\gamma,
\end{equation}
then at least one $a\in\atlas$ is feasible for the target task.
\end{theorem}

\begin{proof}
Assumptions \ref{ass:coordinate} and \ref{ass:support}, followed by the
triangle inequality, place $s_c$ within
$\Delta_{\mathrm{task}}+2\epsilon_\eta$ of $x_c$ in the learned coordinate.
Assumption \ref{ass:atlas-cover} supplies $a\in\atlas$ at an additional
distance at most $r_{\mathrm{cover}}$.  By \eqref{eq:lipschitz} and
\eqref{eq:coverage-condition},
\[
 m_q(a)\le m_q(x_c)+L_m(r_{\mathrm{cover}}+Delta_{\mathrm{task}}
 +2\epsilon_\eta)\le -\gamma+\gamma=0.
\]
Thus $a$ is feasible.
\end{proof}

\begin{corollary}[Nominal-dimension independence]\label{cor:dimension}
If $r_{\mathrm{cover}}$, $\Delta_{\mathrm{task}}$, $\epsilon_\eta$, $L_m$,
and $\gamma$ remain controlled as $d$ changes, the sufficient condition
\eqref{eq:coverage-condition} does not depend on $d$.
\end{corollary}

Corollary \ref{cor:dimension} is a conditional effective-dimension result, not
an assertion that the constants are automatically stable.  The finite-library
post-run audit verifies positive slack in all three synthetic domains (Table
\ref{tab:coverage-audit}) under registered library quantities.  It does not
certify a global Lipschitz constant for every policy, so the global theorem
remains an explicit modeling condition.

\begin{table}[t]
\centering
\caption{Post-run finite-library coverage audit.  Slack is safe radius minus
the audited atlas/support radius.  These diagnostics were not used for target
selection.}
\label{tab:coverage-audit}
\begin{tabular}{lrrrr}
\toprule
Domain & $r_{cover}$ & Support shift & Safe radius & Slack \\
\midrule
FactorShock & 0.0502 & 0.0000 & 0.3277 & 0.2775 \\
Inventory & 0.0307 & 0.0000 & 0.1322 & 0.1015 \\
Queue & 0.0384 & 0.0019 & 0.1298 & 0.0895 \\
\bottomrule
\end{tabular}

\end{table}

\subsection{Fail-Closed Monotone Extension}

Let $U$ and $L$ denote the upper and lower normalized constant policies.  The
source-monotone test of Section \ref{sec:atlas} either replaces one atlas entry
with an admitted endpoint or returns the base atlas unchanged.

\begin{proposition}[Fail-closed identity and budget]\label{prop:failclosed}
If neither source direction satisfies the admission threshold, the extension
returns the base atlas exactly.  In all cases its cardinality is at most $n_0$.
\end{proposition}

\begin{proposition}[Transferred endpoint safety]\label{prop:endpoint}
Suppose all source tasks admit $U$, the target margin is nonincreasing in the
zero-frequency coordinate, and $U$ is coordinate-wise maximal.  If any target
policy is feasible, then $U$ is feasible.  A symmetric statement holds for
the admitted lower endpoint and a nondecreasing target margin.
\end{proposition}

\begin{proof}
For a feasible witness $x_s$, maximality gives
$\bar z(x_s)\le\bar z(U)$.  Target monotonicity gives
$m_q(U)\le m_q(x_s)\le0$.  The lower-endpoint case reverses the order.
\end{proof}

The source rank statistic is only the admission rule.  Target monotonicity is
the additional identifiability condition and is not inferred by logic alone.

\subsection{Optimizer-Independent Terminal Verification}

After search, the backend freezes an ordered shortlist of at most three
policies: an objective challenger, the backend recommendation, and a
safe-interior support policy.  Candidate $j$ is tested with fresh samples and
error allocation $\delta_j$.  Let $E_j$ be the event that candidate $j$ is
certified even though it is unsafe.

\begin{theorem}[Ordered-shortlist safety]\label{thm:verifier}
If $\Prob(E_j)\le\delta_j$ for every frozen candidate, then the probability
that the first certified shortlist member is unsafe is at most
$\sum_j\delta_j$.
\end{theorem}

\begin{proof}
Unsafe deployment implies $E_j$ for at least one candidate that could be
selected.  The union bound yields the result; no property of the optimizer is
used.
\end{proof}

The challenger can replace a certified incumbent only if a one-sided upper
confidence bound on their paired objective difference is negative.  Let
$\delta_{\mathrm{obj}}$ bound failure of that objective upper bound.

\begin{corollary}[Joint deployment and switch bound]\label{cor:joint}
If $\sum_j\delta_j\le\delta_{\mathrm{safe}}$, then the probability of unsafe
deployment or an incorrect objective switch is at most
$\delta_{\mathrm{safe}}+\delta_{\mathrm{obj}}$.
\end{corollary}

The implementation spends familywise error $0.05$ across the frozen shortlist
and uses eight fresh paired objective replications per compared policy when an
objective switch is attempted.  Verification samples never update the search
posterior.

For the external energy experiment, one verification replication is a fresh
calendar window and success means satisfying the reliability requirement in
that window.  If $K\sim\mathrm{Binomial}(n,p)$, the probability of $n$ successes
is $p^n$.

\begin{proposition}[Exact all-success certificate]\label{prop:binomial}
For any infeasible reliability probability $p\le p_{\mathrm{req}}$,
\begin{equation}
 \Prob_p(K=n)=p^n\le p_{\mathrm{req}}^n.
\end{equation}
Error spending over a frozen finite shortlist gives the corresponding
familywise bound.
\end{proposition}

The registered first stage uses 80 independent windows.  At the 0.95 required
probability and registered confidence level, the one-sided Clopper--Pearson
rule certifies only the all-success count; the implementation equivalence is
regression tested and Proposition \ref{prop:binomial} is machine checked from
the binomial law.

\subsection{Optional Cumulative Heteroscedastic Calibration}

Some simulation tasks provide observable local exposures $A(x)$ and shared
shock exposures $N(x)$.  The optional cumulative factor model is
\begin{equation}\label{eq:hvd}
 v_C(x)=A(x)^\top\Lambda A(x)+N(x)^\top B N(x)
       +N(x)^\top\omega+v_0,
\end{equation}
with nonnegative diagonal $\Lambda$, positive-semidefinite $B$, and projected
nonnegative linear/floor terms.  This decomposition recovers local,
shared-shock, linear, and floor contributions and is used only in the HVD
diagnostic.  The primary atlas and verifier do not require it.  The distinction
matters because better variance fit is not the same as better optimization;
Section \ref{sec:hvd-results} reports both outcomes.
